\documentclass[11pt]{article}

% -------------------------
% Provisional Patent Draft (LaTeX)
% Bias-Conditioned Low-Precision Neural Network Training
% -------------------------

\usepackage[a4paper,margin=1in]{geometry}
\usepackage{amsmath,amssymb,amsfonts}
\usepackage{bm}
\usepackage{booktabs}
\usepackage{hyperref}
\usepackage{enumitem}
\usepackage{graphicx}
\usepackage{caption}
\usepackage{float}
\usepackage{listings}
\usepackage{multicol}
\usepackage{xcolor}

\hypersetup{
  colorlinks=true,
  linkcolor=blue,
  urlcolor=blue,
  citecolor=blue
}

\lstset{
  basicstyle=\ttfamily\small,
  breaklines=true,
  frame=single,
  rulecolor=\color{black!30}
}

\title{\bf Provisional Patent Specification:\\
Bias-Conditioned Low-Precision Neural Network Training via Rank-One Mean Decomposition}
\author{Inventor(s): Li Shang \;\; (and others, if applicable)}
\date{\today}

\begin{document}
\maketitle

\vspace{-0.5em}
\noindent\textbf{Applicant (optional):} \underline{\hspace{0.7\linewidth}}\\
\textbf{Correspondence Address (optional):} \underline{\hspace{0.62\linewidth}}\\
\textbf{Email (optional):} \underline{\hspace{0.79\linewidth}}

\begin{abstract}
This disclosure describes systems and methods for stabilizing low-precision training of neural networks by isolating and separately processing a coherent bias component that dominates numerical dynamic range. In representative embodiments, an activation tensor is decomposed into (i) a rank-one mean component and (ii) a residual component. The two components are independently quantized and recombined within matrix multiplications (e.g., GeMM) for the forward pass. A corresponding mean--residual decomposition may be applied symmetrically to gradient tensors in the backward pass, enabling stable low-bit training without performing expensive spectral decompositions such as SVD. The disclosed techniques are hardware-friendly, using reduction operations and elementwise kernels, and can be integrated into accelerator training pipelines (e.g., GPU tensor cores, ASICs). Optional extensions include residual low-rank spike conditioning, blockwise/groupwise mean decomposition, mixed-precision encoding of the mean component, and distributed implementations. The disclosed methods mitigate outlier-driven quantization scale hijacking and reduce second-moment inflation in adaptive optimization, improving training stability and efficiency.
\end{abstract}

\tableofcontents
\newpage

% ============================================================
\section{Field of the Invention}
The present disclosure relates to machine learning systems and neural network training. More particularly, it relates to methods and systems for stabilizing reduced-precision (e.g., FP8, FP4, INT8-like, block floating point) training and/or inference by decomposing tensors into coherent bias components and residual components and applying independent quantization and computation pipelines.

% ============================================================
\section{Background}
\subsection{Reduced-precision training and instability}
Neural network training increasingly relies on reduced-precision arithmetic to improve throughput, reduce memory bandwidth pressure, and lower energy consumption. However, low-precision formats impose limited representable dynamic range and coarse quantization bins. As a result, extreme tensor values (outliers) can dominate quantization scaling, compressing informative residual variation and causing training instability, divergence, or degraded convergence.

\subsection{Existing mitigation approaches}
Existing approaches often attempt to suppress outliers or anisotropy using one or more of the following:
\begin{itemize}[leftmargin=1.2em]
\item Explicit spectral decompositions (e.g., SVD-based low-rank suppression or whitening);
\item Orthogonal rotations or structured transforms;
\item Rescaling or redistribution methods that shift magnitude between weights and activations;
\item Clipping, outlier suppression, and other heuristic stabilizers.
\end{itemize}
While effective in some settings, spectral decompositions and iterative orthogonalization can be computationally expensive, memory intensive, and poorly aligned with accelerator hardware pipelines.

\subsection{Need for hardware-friendly alternatives}
There remains a need for methods that provide comparable stability benefits while being implementable using inexpensive operations such as reductions and elementwise kernels and without requiring repeated spectral decompositions during training.

% ============================================================
\section{Summary of the Invention}
This disclosure provides methods and systems that:
\begin{enumerate}[leftmargin=1.3em]
\item compute a coherent bias component (e.g., feature-wise mean) of a tensor;
\item represent said bias component as a rank-one (or low-rank) structure;
\item compute a residual tensor by subtracting the bias component;
\item independently quantize the bias component and residual tensor (possibly with different quantizers and scales);
\item perform neural network computations (e.g., matrix multiplications) using the independently quantized components and recombine their contributions;
\item optionally apply a symmetric decomposition and independent quantization to backward-pass gradient tensors;
\item optionally apply lightweight conditioning to the residual (e.g., low-rank spike conditioning/whitening) without full-layer spectral decomposition.
\end{enumerate}
In representative embodiments, the coherent bias is the feature-wise mean vector, yielding a rank-one bias matrix. This rank-one mean component is quantized separately, preventing it from hijacking quantization scale intended for the residual semantic signal. The result is improved stability and performance in low-bit training, with minimal overhead.

% ============================================================
\section{Brief Description of the Drawings (Optional)}
The following drawings are illustrative and may be included in a later non-provisional filing; they are not required for a provisional but are recommended.
\begin{itemize}[leftmargin=1.2em]
\item \textbf{FIG. 1}: System diagram of a training pipeline with bias-conditioned quantization modules integrated with accelerator kernels.
\item \textbf{FIG. 2}: Flowchart of forward-pass mean--residual decomposition, independent quantization, and recomposed GeMM.
\item \textbf{FIG. 3}: Flowchart of backward-pass symmetric decomposition for gradient tensors and computation of $\partial L/\partial X$ and $\partial L/\partial W$.
\item \textbf{FIG. 4}: Blockwise/groupwise mean decomposition and per-block quantization scale computation.
\item \textbf{FIG. 5}: Optional residual spike conditioning applied selectively to the residual component.
\end{itemize}

% ============================================================
\section{Definitions and Notation}
For clarity, the following terms are used throughout:

\paragraph{Activation tensor.}
A tensor (or matrix after reshaping) representing intermediate forward-pass values of a layer.

\paragraph{Gradient tensor.}
A tensor (or matrix) representing backward-pass gradients with respect to a layer output.

\paragraph{Feature-wise mean vector.}
A vector of means computed across a token/batch dimension for each feature/channel dimension.

\paragraph{Rank-one mean tensor (bias tensor).}
A tensor constructed as an outer product of an all-ones vector with a mean vector, yielding a rank-one structure in the matrix view.

\paragraph{Residual tensor.}
A tensor defined as activation minus the rank-one mean tensor.

\paragraph{Independent quantization.}
Quantizing the bias tensor and residual tensor separately, optionally using distinct scaling rules, block partitions, bit-widths, and/or stochastic rounding.

\paragraph{Low-precision format.}
Any reduced-precision numerical representation including FP8, FP4, block floating point, INT8-like formats, or custom formats with reduced mantissa/exponent.

\paragraph{GeMM.}
General matrix-matrix multiplication; in practice may be implemented via specialized accelerator kernels (tensor cores, systolic arrays, etc.).

\paragraph{Token/batch dimension.}
A dimension representing batch size and/or sequence length (e.g., $l=b\cdot s$).

% ============================================================
\section{Detailed Description of Embodiments}

\subsection{Overview}
The disclosed methods operate on tensors during training and/or inference. For training, operations occur in the forward pass (activations) and optionally in the backward pass (gradients). The key insight is that a coherent bias component (often rank-one when represented as a mean outer product) can dominate extreme values and quantization scaling. By separating this component and quantizing it independently from the residual, the residual retains improved effective resolution in low precision.

\subsection{Representative forward-pass embodiment (mean--residual splitting)}
Consider a linear layer with weight matrix $W\in\mathbb{R}^{m\times n}$. Let an input activation tensor be reshaped into a matrix
\[
X \in \mathbb{R}^{l\times m},
\]
where $l$ is the flattened token/batch dimension. The exact reshape is implementation-dependent and may correspond to batch size $\times$ sequence length.

\paragraph{Compute feature-wise mean.}
Compute the mean vector
\[
\mu_X \;=\; \frac{1}{l}\bm{1}^\top X \in \mathbb{R}^{m},
\]
where $\bm{1}\in\mathbb{R}^{l}$ is an all-ones vector.

\paragraph{Construct rank-one mean tensor.}
Construct the mean (bias) matrix
\[
M_X \;=\; \bm{1}\mu_X^\top \in \mathbb{R}^{l\times m}.
\]

\paragraph{Compute residual.}
Compute the residual
\[
X_R \;=\; X - M_X.
\]

\paragraph{Independent quantization.}
Let $Q_b(\cdot)$ denote a quantizer into a reduced-precision format (e.g., block floating point FP4/FP8).
Quantize components independently:
\[
\overline{M}_X = Q_b(M_X), \qquad \overline{X}_R = Q_b(X_R), \qquad \overline{W}=Q_b(W).
\]
In some embodiments, $\overline{M}_X$ may use a different precision and/or scaling policy than $\overline{X}_R$.

\paragraph{Recomposed computation.}
Compute the output
\[
\widehat{Y} \;=\; (\overline{M}_X + \overline{X}_R)\,\overline{W}.
\]
In some embodiments, the mean matrix $M_X$ is never explicitly materialized; rather, $\mu_X$ is computed, and subtraction is performed on-the-fly to form $X_R$.

\subsection{Representative backward-pass embodiment (symmetric decomposition)}
Let $D \in \mathbb{R}^{l\times n}$ denote the gradient matrix with respect to the pre-activation output $Y$ (after reshaping),
\[
D \;=\; \frac{\partial L}{\partial Y}.
\]
Compute mean and residual for $D$:
\[
\mu_D = \frac{1}{l}\bm{1}^\top D \in \mathbb{R}^{n},\qquad
M_D = \bm{1}\mu_D^\top,\qquad
D_R = D - M_D.
\]
Quantize independently:
\[
\overline{M}_D = Q_b(M_D),\qquad \overline{D}_R=Q_b(D_R).
\]

\paragraph{Compute gradients.}
Compute quantized approximations of gradients:
\[
\widehat{\frac{\partial L}{\partial X}} \;=\; (\overline{M}_D + \overline{D}_R)\,\overline{W}^\top,
\]
\[
\widehat{\frac{\partial L}{\partial W}} \;=\; (\overline{M}_X + \overline{X}_R)^\top(\overline{M}_D + \overline{D}_R).
\]
In some embodiments, $\widehat{\partial L/\partial W}$ is formed with mixed-precision accumulation and/or blockwise reduction.

\subsection{Quantizer embodiments}
The quantizer $Q_b(\cdot)$ may be implemented as:
\begin{itemize}[leftmargin=1.2em]
\item Blockwise floating point (e.g., FP8/FP4 with per-block scale);
\item Per-channel or per-group quantization;
\item Stochastic rounding quantization;
\item Learned scale quantization;
\item Custom exponent/mantissa configurations.
\end{itemize}
In some embodiments, separate scales are computed for bias and residual:
\[
s_M = \mathrm{Scale}(M_X),\qquad s_R=\mathrm{Scale}(X_R).
\]
The scales may be computed from max-abs, percentile statistics, or robust estimators.

\subsection{Hardware-friendly implementation}
The disclosed method may be implemented using only:
\begin{itemize}[leftmargin=1.2em]
\item Reduction kernels for computing $\mu_X$ and/or $\mu_D$,
\item Elementwise subtraction kernels for forming residuals,
\item Existing quantized GeMM kernels for matrix multiplications,
\item Optional fused kernels that combine mean computation, subtraction, and quantization.
\end{itemize}
In contrast to SVD-based methods, no eigen decomposition, QR decomposition, or iterative orthogonalization is required.

% ============================================================
\section{Optional Extensions and Variations (Defensive Breadth)}

\subsection{Partial mean scaling}
Instead of fully removing the mean tensor, a scaling factor $\alpha\in[0,1]$ may be applied:
\[
X_R = X - \alpha M_X,\qquad D_R = D - \alpha M_D.
\]
The factor $\alpha$ may depend on layer index, training step, token type, or measured outlier statistics.

\subsection{Groupwise / blockwise mean decomposition}
Mean vectors may be computed:
\begin{itemize}[leftmargin=1.2em]
\item Per group of channels/features (group size $g$),
\item Per attention head (transformer),
\item Per quantization block/tile aligned to hardware layouts.
\end{itemize}
For example, partition features into groups $G_k$ and compute $\mu_X^{(k)}$ for each group.

\subsection{Mixed precision for bias vs residual}
In some embodiments, the bias component is encoded at higher precision than the residual:
\[
\overline{M}_X = Q_{b_1}(M_X),\qquad \overline{X}_R = Q_{b_2}(X_R),\quad b_1>b_2.
\]
This preserves coherent bias while maximizing compression of residual variation.

\subsection{Selective application across layers}
The method may be applied:
\begin{itemize}[leftmargin=1.2em]
\item only at the embedding layer,
\item only at early layers,
\item only at layers detected to have large coherent bias,
\item only during certain training phases.
\end{itemize}

\subsection{Residual spike conditioning (optional lightweight whitening)}
In addition to removing a rank-one mean, residual low-rank spikes may be conditioned. In some embodiments, apply a lightweight low-rank operation to $X_R$ (or $D_R$) such as:
\begin{itemize}[leftmargin=1.2em]
\item Approximate whitening on a small number of directions,
\item Randomized low-rank approximation,
\item Incremental PCA/Oja-style updates to track top directions,
\item Hadamard/rotation transforms applied to residual only.
\end{itemize}
Such conditioning is applied selectively and does not require full per-layer SVD.

\subsection{Distributed training embodiments}
In distributed data-parallel training, mean computation may be:
\begin{itemize}[leftmargin=1.2em]
\item local per device,
\item aggregated across devices via all-reduce to form a global mean,
\item performed per microbatch to reduce latency.
\end{itemize}
In some embodiments, the mean vector is computed prior to gradient synchronization to reduce outlier effects in communication and optimizer states.

\subsection{Inference-only embodiments}
Although described for training, the method may be used in inference-time quantization:
\begin{itemize}[leftmargin=1.2em]
\item compute mean/residual per token block,
\item quantize separately,
\item improve low-bit inference quality and stability.
\end{itemize}

% ============================================================
\section{Optimizer Integration (Optional Embodiments)}
Adaptive optimizers maintain second-moment statistics (e.g., Adam/AdamW). Because gradients of linear layers have outer-product structure $g=\delta x^\top$, squared activation magnitudes can inflate second moments. By reducing coherent bias in activations, the disclosed method may reduce second-moment inflation and improve effective learning-rate balance. In some embodiments, mean removal is applied upstream of optimizer updates and/or selectively enabled when optimizer statistics exceed thresholds.

% ============================================================
\section{Advantages}
Representative advantages include:
\begin{itemize}[leftmargin=1.2em]
\item Reduced outlier-driven quantization scale hijacking;
\item Improved effective resolution of residual (semantic) signal in low-bit formats;
\item Improved numerical stability and convergence in low-bit training;
\item Reduction in extreme second-moment statistics in adaptive optimizers (optional);
\item Hardware-friendly implementation using reductions and elementwise operations;
\item Avoidance of expensive spectral decompositions (e.g., SVD).
\end{itemize}

% ============================================================
\section{Example Implementation (Pseudo-Workflow)}
\begin{lstlisting}[language=,caption={Representative forward/backward bias-conditioned training workflow (pseudo-code)}]
# Forward
X = reshape(activation_tensor, [l, m])
muX = mean(X, axis=0)               # reduction
XR  = X - ones(l,1) @ muX[None,:]   # residual (rank-one subtraction)
MX  = ones(l,1) @ muX[None,:]       # implicit or explicit

MX_q = quantize(MX, scheme_M)
XR_q = quantize(XR, scheme_R)
W_q  = quantize(W,  scheme_W)

Yhat = (MX_q + XR_q) @ W_q

# Backward (optional)
D = reshape(dL_dY, [l, n])
muD = mean(D, axis=0)
DR  = D - ones(l,1) @ muD[None,:]
MD  = ones(l,1) @ muD[None,:]

MD_q = quantize(MD, scheme_MD)
DR_q = quantize(DR, scheme_DR)

dXhat = (MD_q + DR_q) @ W_q.T
dWhat = (MX_q + XR_q).T @ (MD_q + DR_q)
\end{lstlisting}

% ============================================================
\section{Draft Claim Set (Optional for Provisional)}
\subsection*{Note}
Claims are not required for a provisional application, but including a draft claim set can strengthen clarity and later conversion.

\subsection{Independent Claims (Method)}
\begin{enumerate}[leftmargin=1.3em]
\item A computer-implemented method for training a neural network using low-precision numerical representation, comprising:
(a) receiving an activation tensor associated with a layer of the neural network;
(b) computing a feature-wise mean vector from the activation tensor by a reduction operation across a token or batch dimension;
(c) constructing a rank-one mean tensor from the feature-wise mean vector;
(d) generating a residual activation tensor by subtracting the rank-one mean tensor from the activation tensor;
(e) independently quantizing the rank-one mean tensor and the residual activation tensor into a reduced numerical precision format;
(f) performing a matrix multiplication between the independently quantized tensors and a quantized weight tensor to generate an output tensor; and
(g) reconstructing the output tensor from contributions of the independently quantized rank-one mean tensor and the independently quantized residual activation tensor.

\item The method of claim 1, further comprising:
(a) receiving a gradient tensor corresponding to the output tensor;
(b) computing a feature-wise mean of the gradient tensor;
(c) constructing a gradient rank-one mean tensor;
(d) generating a residual gradient tensor by subtracting the gradient rank-one mean tensor from the gradient tensor;
(e) independently quantizing the gradient rank-one mean tensor and the residual gradient tensor; and
(f) computing gradients with respect to weights and input activations using matrix multiplications involving the independently quantized gradient components.
\end{enumerate}

\subsection{Independent Claims (System / CRM)}
\begin{enumerate}[leftmargin=1.3em, resume]
\item A neural network training system comprising one or more processors and memory storing instructions that, when executed, cause the one or more processors to perform the method of claim 1.

\item A non-transitory computer-readable medium storing instructions that, when executed by one or more processors, cause the one or more processors to perform the method of claim 1.
\end{enumerate}

\subsection{Dependent Claims (Examples)}
\begin{enumerate}[leftmargin=1.3em, resume]
\item The method of claim 1, wherein the reduced numerical precision format comprises a 4-bit floating point format.
\item The method of claim 1, wherein quantization scales for the rank-one mean tensor and the residual activation tensor are computed independently.
\item The method of claim 2, wherein the layer is a transformer layer and the mean vector is computed per attention head.
\item The method of claim 1, wherein constructing the rank-one mean tensor is performed implicitly without materializing the rank-one mean tensor in memory.
\item The method of claim 1, further comprising applying a partial mean scaling factor $\alpha\in[0,1]$ when generating the residual activation tensor.
\item The method of claim 1, wherein the method is performed without computing singular value decomposition, eigen decomposition, or QR decomposition.
\item The method of claim 1, wherein the rank-one mean tensor is quantized using a higher precision than the residual activation tensor.
\item The method of claim 1, wherein the feature-wise mean vector is computed per quantization block aligned to hardware tile boundaries.
\item The method of claim 2, further comprising applying residual low-rank conditioning to the residual activation tensor without performing full-layer spectral decomposition.
\end{enumerate}

% ============================================================
\section{Appendix (Optional): Embodiment Checklists}
\subsection*{A. Checklist of covered variations}
\begin{multicols}{2}
\begin{itemize}[leftmargin=1.2em]
\item Full mean subtraction ($\alpha=1$)
\item Partial mean subtraction ($\alpha<1$)
\item Groupwise mean
\item Per-head mean
\item Per-block mean
\item Forward-only
\item Forward + backward symmetric
\item Same-precision bias/residual
\item Mixed-precision bias/residual
\item Stochastic rounding
\item Blockwise scaling
\item Distributed all-reduce mean
\item Inference-only embodiment
\item Optional residual spike conditioning
\end{itemize}
\end{multicols}

\end{document}